\documentclass[
	fontsize=12pt,          % default font size 12pt
	paper=a4,               % DIN A4 page format
	numbers=noenddot,       % remove dots behind chapter numbers (e.g. 1.5 not 1.5.)
	listof=totoc,           % add list of figures, tables, etc. to ToC
	listof=entryprefix,     % add entry name to figures, tables, etc.
	listof=nochaptergap,    % no chapter gap for figures, tables, etc.
	bibliography=totoc,     % add bibliography to ToC but without a chapter number
	parskip=half            % half line spacing between paragraphs
	openany                 % chapters can start on any page
]{scrbook}
\usepackage{tcolorbox}

\include{../../for_latex/preamble}
\title{\Huge\textbf{Kinoreservierung}}
\author{Luis Staudt}
\date{}

% Code-Highlighting für Java
\definecolor{codegray}{rgb}{0.5,0.5,0.5}
\definecolor{codepurple}{rgb}{0.58,0,0.82}
\definecolor{backcolour}{rgb}{0.95,0.95,0.92}
\definecolor{keywordcolor}{rgb}{0.8,0.47,0.13}

\lstdefinestyle{java}{
	language=,
	basicstyle=\ttfamily\small,
	keywordstyle=\bfseries\color{blue!70!black},
	stringstyle=\color{red!60!black},
	commentstyle=\itshape\color{gray},
	numbers=left,
	numberstyle=\tiny\color{gray},
	numbersep=8pt,
	frame=single,
	framerule=0.4pt,
	rulecolor=\color{gray!50},
	breaklines=true,
	showstringspaces=false,
	tabsize=4,
	xleftmargin=1.5em,
	framexleftmargin=1.5em,
	aboveskip=0.8em,
	belowskip=0.6em
}
\lstset{style=java}

\newtcolorbox{hinweis}[1][Hinweis]{
	colback=blue!5,
	colframe=blue!40!black,
	fonttitle=\bfseries,
	title=#1,
	boxrule=0.5pt,
	arc=2pt
}

\newtcolorbox{beispiel}[1][Beispiel]{
	colback=green!5,
	colframe=green!40!black,
	fonttitle=\bfseries,
	title=#1,
	boxrule=0.5pt,
	arc=2pt
}

\begin{document}

% --- Titelblock ---
\begin{center}
	\rule{\textwidth}{0.8pt}\\[0.6em]
	{\LARGE\bfseries Intensivtutorium}\\[0.3em]
	{\large Programmierung}\\[0.6em]
	\begin{tabular}{r l @{\hspace{2cm}} r l}
		\textbf{Semester:}      & Sommersemester 2026                    & \textbf{Tutor:}         & Luis Staudt \\
		\textbf{Studiengänge:}  & MIB, OMB, PEB, MVB         & & \\
	\end{tabular}\\[0.6em]
	\rule{\textwidth}{0.8pt}
\end{center}
\vspace{1em}

% --- Seitenkopf und -fuß (ab Seite 2) ---
\pagestyle{fancy}
\fancyhf{}
\lhead{\small Intensivtutorium -- Grundlagen der Programmierung}
\rhead{\small SS 2026}
\cfoot{\thepage}
\renewcommand{\headrulewidth}{0.4pt}


% =============================================
% Aufgabe 1 -- Kinoreservierung
% =============================================

\section*{Aufgabe 1 -- Kinoreservierung}

Ein kleines Kino verwaltet seine Sitzplätze mit einem zweidimensionalen Array. Jeder Eintrag steht für einen Sitzplatz: \texttt{true} bedeutet \emph{belegt}, \texttt{false} bedeutet \emph{frei}.

\begin{lstlisting}
boolean[][] kino = new boolean[reihen][plaetze];
\end{lstlisting}

Der Saal hat \texttt{reihen} Reihen mit jeweils \texttt{plaetze} Sitzen. Zu Beginn sind alle Plätze frei.

\begin{center}
	\begin{tikzpicture}[
		frei/.style={draw, minimum size=0.55cm, fill=green!15, font=\footnotesize\ttfamily},
		belegt/.style={draw, minimum size=0.55cm, fill=red!25, font=\footnotesize\ttfamily},
	]
		% Reihe 0
		\node[anchor=east, font=\footnotesize] at (-0.2, 0) {Reihe 0:};
		\foreach \p/\s in {0/frei,1/frei,2/belegt,3/frei,4/frei,5/frei,6/frei,7/frei} {
			\node[\s] at (\p*0.65, 0) {\ifx\s\belegt X\else O\fi};
		}
		% helper macro workaround
		\node[belegt] at (2*0.65, 0) {X};
		
		% Reihe 1
		\node[anchor=east, font=\footnotesize] at (-0.2, -0.75) {Reihe 1:};
		\foreach \p in {0,...,7} {
			\node[frei] at (\p*0.65, -0.75) {O};
		}
		
		% Reihe 2
		\node[anchor=east, font=\footnotesize] at (-0.2, -1.5) {Reihe 2:};
		\node[frei]   at (0*0.65, -1.5) {O};
		\node[belegt] at (1*0.65, -1.5) {X};
		\node[belegt] at (2*0.65, -1.5) {X};
		\node[belegt] at (3*0.65, -1.5) {X};
		\foreach \p in {4,...,7} {
			\node[frei] at (\p*0.65, -1.5) {O};
		}
		
		% Reihe 3
		\node[anchor=east, font=\footnotesize] at (-0.2, -2.25) {Reihe 3:};
		\foreach \p in {0,...,7} {
			\node[frei] at (\p*0.65, -2.25) {O};
		}
		
		% Reihe 4
		\node[anchor=east, font=\footnotesize] at (-0.2, -3.0) {Reihe 4:};
		\foreach \p in {0,...,4} {
			\node[frei] at (\p*0.65, -3.0) {O};
		}
		\node[belegt] at (5*0.65, -3.0) {X};
		\node[belegt] at (6*0.65, -3.0) {X};
		\node[frei]   at (7*0.65, -3.0) {O};
		
		% Spaltenindizes
		\foreach \p in {0,...,7} {
			\node[font=\tiny\color{gray}] at (\p*0.65, 0.55) {\p};
		}
		
		% Legende
		\node[frei, anchor=west] at (6.2, -0.5) {O};
		\node[anchor=west, font=\footnotesize] at (6.85, -0.5) {= frei};
		\node[belegt, anchor=west] at (6.2, -1.15) {X};
		\node[anchor=west, font=\footnotesize] at (6.85, -1.15) {= belegt};
		
		% Index-Hinweis
		\draw[-{Stealth}, gray, thick] (-0.5, 0.55) -- (-0.5, -3.0)
		node[midway, left, font=\tiny\color{gray}] {r};
		\draw[-{Stealth}, gray, thick] (-0.15, 0.85) -- (7*0.65+0.25, 0.85)
		node[midway, above, font=\tiny\color{gray}] {p};
	\end{tikzpicture}
\end{center}

\begin{hinweis}
	Der Index \texttt{kino[r][p]} steht für Reihe~\texttt{r}, Platz~\texttt{p} (jeweils ab~0).
\end{hinweis}

\bigskip
Implementiere die folgenden Methoden:

\begin{enumerate}[label=\textbf{\alph*)}]

\item \textbf{Saal anzeigen} -- Gibt den Saal zeilenweise auf der Konsole aus. Verwende \texttt{X} für belegte und \texttt{O} für freie Plätze. Vor jeder Zeile soll die Reihennummer stehen.

\begin{lstlisting}
static void anzeigen(boolean[][] kino)
\end{lstlisting}

\item \textbf{Platz reservieren} -- Reserviert den Platz in Reihe~\texttt{r}, Position~\texttt{p}. Ist der Platz bereits belegt oder sind die Indizes ungültig, soll eine Fehlermeldung auf der Konsole ausgegeben und \texttt{false} zurückgegeben werden. Bei erfolgreicher Reservierung wird \texttt{true} zurückgegeben.

\begin{lstlisting}
static boolean reservieren(boolean[][] kino, int r, int p)
\end{lstlisting}

\item \textbf{Freie Plätze zählen} -- Gibt die Gesamtanzahl aller freien Plätze im Kino zurück.

\begin{lstlisting}
static int freiePlaetze(boolean[][] kino)
\end{lstlisting}

\item \textbf{Reihe komplett frei?} -- Prüft, ob in der angegebenen Reihe \emph{alle} Plätze frei sind.

\begin{lstlisting}
static boolean reiheFrei(boolean[][] kino, int reihe)
\end{lstlisting}

\item \textbf{Gruppenreservierung} -- Eine Gruppe von \texttt{anzahl} Personen möchte nebeneinander in \emph{derselben Reihe} sitzen. Die Methode sucht den ersten zusammenhängenden freien Bereich, der groß genug ist, reserviert die Plätze und gibt die Reihe zurück. Wird kein passender Bereich gefunden, wird \texttt{-1} zurückgegeben.

\begin{lstlisting}
static int gruppenReservierung(boolean[][] kino, int anzahl)
\end{lstlisting}

\begin{hinweis}
	Für Teilaufgabe e) muss in jeder Reihe nach einem zusammenhängenden Block aus \texttt{anzahl} freien Plätzen gesucht werden. Überlege dir, wie du die Länge des aktuellen freien Bereichs beim Durchlaufen mitzählen kannst.
\end{hinweis}

\item \textbf{Belegungsgrad} -- Berechnet den prozentualen Anteil der belegten Plätze und gibt ihn als \texttt{double} zurück (z.\,B.\ 42.5 für 42{,}5\,\%).

\begin{lstlisting}
static double belegungsgrad(boolean[][] kino)
\end{lstlisting}

\end{enumerate}

\end{document}
